\section{Introduction}
\label{sec:intro}

Data spaces are emerging as a
societal, legal and technical architectural paradigm
at the centre of the European Commission's policy agenda
for enabling trusted data sharing across organisational
and sectoral boundaries~\cite{farrell_european_2023}.
Their objective is to allow participants
to exchange and reuse data while maintaining control over
how that data is accessed and processed.
Achieving this balance between
data sharing and data sovereignty
requires usage control mechanisms
that can express and enforce the conditions
under which data may be used~\cite{eitel_usage_2021}.
These conditions can be derived not only from technical
but also from regulatory and contractual requirements.
In this context,
the Data Spaces Support Centre (DSSC) blueprint~\cite{data_spaces_support_centre_summary_2026}
highlights the importance of integrating legal and technical capabilities,
specifically addressing \emph{regulatory compliance and contractual frameworks}
alongside \emph{access and usage policies enforcement}.
These building blocks together aim to ensure that
the rights and obligations governing data sharing can be
defined, negotiated, and enforced in an interoperable manner.
Agreements between data providers and data consumers
must therefore be represented in a way that
supports automated interpretation, negotiation,
and enforcement during the lifecycle of data transactions.

A central challenge in this context
is the representation of key agreements
while preserving their legal meaning.
Data sharing policies must capture
permissions, prohibitions, and obligations
governing data access and usage,
while also reflecting constraints
derived from sectoral rules, 
contract law and intellectual property rights,
or data-related legislation.
At the same time,
these policies must be semantically interoperable
to support the federated nature of data spaces,
ensuring ``that data can be accurately and
consistently interpreted and used by
different people and systems''~\cite{steinbuss_semantic_2025}.
Two complementary semantic specifications provide
a promising foundation for addressing this challenge.
The W3C Open Digital Rights Language (ODRL) standard~\cite{iannella_odrl_2018}
provides a general-purpose policy model
for expressing rules
governing the use of digital assets,
including permissions, prohibitions, and obligations
which can be further restricted with contextual constraints.
Furthermore,
ODRL is recommended in the DSSC blueprint
as a standard representation for machine-readable
access and usage policies within data spaces.
While ODRL provides the structural framework for policy representation,
it intentionally remains domain- and jurisdiction-agnostic and
does not include explicit concepts related to legal compliance.
The Data Privacy Vocabulary (DPV)~\cite{pandit_data_2024},
built as a `community standard',
complements ODRL by providing a rich semantic vocabulary
containing taxonomies to describe legal and sectorial requirements
related to data and AI regulation.
Combining these two standards
makes it possible to bridge the gap between
legally grounded governance requirements and
machine-readable policy enforcement mechanisms.
In particular, the use of DPV concepts allows ODRL policies
and documentation to be aligned with the terminologies 
of relevant laws and simplifies the process
of assessing and enforcing compliance.
This allows data usage policies to express legally meaningful constraints
while remaining compatible with existing policy evaluation and enforcement engines.
This approach aligns closely with
the DSSC blueprint's vision of interoperable data spaces
where legal agreements and
technical enforcement mechanisms
are coherently connected.

In this paper,
we argue that ODRL and DPV together
provide a suitable semantic foundation
for legally-aligned policy expression in data spaces.
We propose treating ODRL as the core policy framework
for representing permissions, prohibitions, and
obligations governing data usage,
in line with the DSSC and given the mature status
of ODRL as a W3C standard.
We then propose using DPV as the semantic layer
for modelling legal and governance concepts
referenced within those policies, based on its
wide coverage of EU's legal framework, extensibility
for jurisdictional and sectorial applicability, as well
as a lack of rigid semantics that allows wider reuse.
To support this integration,
we present a structured `mapping' between DPV and ODRL
that aligns DPV concepts with ODRL constructs,
using RDF and SKOS properties,
which allows the integration and reuse of both standards
in a cohesive and semantically meaningful manner.
We also provide the mapping in a machine-readable form
to facilitate implementation,
in accordance with ODRL profiling best practices~\cite{steidl_odrl_2023}.
In addition,
we provide a guidance document for the use of
DPV terms within ODRL policies,
including examples of policies that use the developed mapping.
Together,
these contributions enable data space participants
to express governance rules that are both legally meaningful
and technically enforceable,
thereby supporting scalable and
interoperable policy management
through automation
across data space infrastructures.

In this context,
our work addresses the following challenges
documented by the W3C Dataspaces CG\footnote{https://w3c-cg.github.io/dataspaces/}:
\begin{inlinelist}
    \item[(i)] \textit{Blueprints for actors, roles and responsibilities}---our work provides the means to declare and use these concepts, with legal relevance, through existing standards;
    \item[(ii)] \textit{Interoperable Policy Engines}---our work facilitates improvements to policy engines by providing a legal semantics layer, and also enables legal interoperability across data spaces through shared vocabularies; and 
    \item[(iii)] \textit{Schema Alignment and Semantic Data Transformation}---same as (ii), our work enables creation of schemas and vocabularies within a data space or instance, which can be reused elsewhere based on extending common interpretations through ODRL as a standard for policies and DPV as a standard for expression of legal concepts
\end{inlinelist}.

The remainder of this article is structured as follows:
Section~\ref{sec:background} provides an overview of applicable regulatory requirements and existing work for data protection-aligned policy modelling.
Section~\ref{sec:mapping} describes the developed mapping and its implementation as an ODRL profile.
% In Section~\ref{sec:policies}, we define a specification for ...
In Section~\ref{sec:guidelines}, we provide examples of the application of the profile and discuss the implementation of templates to validate profile-based constraints.
Finally, Section~\ref{sec:conclusions} concludes the paper and provides pointers to future research.